\chapter{Spectral Theorem}

\noindent\textit{Source: Lay \& McDonald, \textit{Linear Algebra and Its Applications}, \S7.1.}

\begin{remark}[Suggested Problems 4]
\textbf{Quiz:} Lay \S7.1.13--22 (orthogonal diagonalization). \textbf{Exam:} Lay \S7.1.25--38 (true/false and proofs).
\end{remark}

%% ================================================================
\section{Symmetric Matrices and Orthogonal Diagonalization}
%% ================================================================

\begin{definition}[Symmetric Matrix]
A matrix $A$ is \textbf{symmetric} if $A^T=A$. Symmetric matrices are square.
\end{definition}

\begin{theorem}[Eigenvectors of Symmetric Matrices]
If $A$ is symmetric, then any two eigenvectors from different eigenspaces are orthogonal.
\end{theorem}

\begin{definition}[Orthogonally Diagonalizable]
An $n\times n$ matrix $A$ is \textbf{orthogonally diagonalizable} if there exist an orthogonal matrix $P$ and a diagonal matrix $D$ such that $A=PDP^T=PDP^{-1}$.
\end{definition}

\begin{theorem}[Spectral Theorem]
An $n\times n$ matrix $A$ is orthogonally diagonalizable \emph{if and only if} $A$ is symmetric. Equivalently: $A$ has $n$ real eigenvalues (counting multiplicities); the dimension of each eigenspace equals the multiplicity of the corresponding eigenvalue; eigenspaces for different eigenvalues are mutually orthogonal.
\end{theorem}

%% ================================================================
\section{Spectral Decomposition}
%% ================================================================

If $A=PDP^T$ with orthonormal eigenvectors $u_1,\ldots,u_n$ and eigenvalues $\lambda_1,\ldots,\lambda_n$, then
\[
A = \lambda_1 u_1 u_1^T + \lambda_2 u_2 u_2^T + \cdots + \lambda_n u_n u_n^T.
\]
Each term $\lambda_i u_i u_i^T$ is a rank-$1$ matrix; $u_i u_i^T$ is the orthogonal projection onto $\operatorname{span}\{u_i\}$.

%% ================================================================
\section{Repeated Eigenvalues}
%% ================================================================

When an eigenvalue has multiplicity $>1$, the eigenspace has dimension equal to the multiplicity. Use Gram--Schmidt on a basis of the eigenspace to obtain orthonormal eigenvectors for the orthogonal diagonalization.

%% ================================================================
\section{How to Orthogonally Diagonalize a Symmetric Matrix}
%% ================================================================

For a symmetric matrix $A$, the workflow is always the same:
\begin{enumerate}
  \item Compute the eigenvalues of $A$.
  \item For each eigenvalue $\lambda$, find a basis of $\operatorname{Nul}(A-\lambda I)$.
  \item If an eigenspace has dimension greater than $1$, apply Gram--Schmidt inside that eigenspace to obtain an orthonormal basis.
  \item Normalize all eigenvectors and place them as columns of $P$.
  \item Put the corresponding eigenvalues on the diagonal of $D$.
\end{enumerate}
Then $A=PDP^T$ with $P^TP=I$.

\begin{remark}
The two checks that matter most in problems are:
\[
P^TP=I
\qquad \text{and} \qquad
A=PDP^T.
\]
If both hold, the orthogonal diagonalization is correct.
\end{remark}

%% ================================================================
\section{Worked Example}
%% ================================================================

Orthogonally diagonalize
\[
A=\begin{pmatrix}
7 & 2\\
2 & 4
\end{pmatrix}.
\]

Since $A^T=A$, the Spectral Theorem applies.

\textbf{Step 1: Find the eigenvalues.}
\[
\det(A-\lambda I)
=
\begin{vmatrix}
7-\lambda & 2\\
2 & 4-\lambda
\end{vmatrix}
=(7-\lambda)(4-\lambda)-4
=\lambda^2-11\lambda+24.
\]
So
\[
\lambda^2-11\lambda+24=(\lambda-8)(\lambda-3),
\]
hence the eigenvalues are $\lambda_1=8$ and $\lambda_2=3$.

\textbf{Step 2: Find orthonormal eigenvectors.}

For $\lambda_1=8$,
\[
A-8I=\begin{pmatrix}-1&2\\2&-4\end{pmatrix},
\]
so $-x+2y=0$, hence $x=2y$. One eigenvector is
\[
v_1=\begin{pmatrix}2\\1\end{pmatrix}.
\]
Normalize:
\[
u_1=\frac{1}{\sqrt{5}}\begin{pmatrix}2\\1\end{pmatrix}.
\]

For $\lambda_2=3$,
\[
A-3I=\begin{pmatrix}4&2\\2&1\end{pmatrix},
\]
so $2x+y=0$, hence $y=-2x$. One eigenvector is
\[
v_2=\begin{pmatrix}1\\-2\end{pmatrix}.
\]
Normalize:
\[
u_2=\frac{1}{\sqrt{5}}\begin{pmatrix}1\\-2\end{pmatrix}.
\]

\textbf{Step 3: Build $P$ and $D$.}
\[
P=\begin{pmatrix}
2/\sqrt{5} & 1/\sqrt{5}\\
1/\sqrt{5} & -2/\sqrt{5}
\end{pmatrix},
\qquad
D=\begin{pmatrix}
8 & 0\\
0 & 3
\end{pmatrix}.
\]

Therefore
\[
\boxed{
A=PDP^T.
}
\]

%% ================================================================
\section{Why This Matters for SVD}
%% ================================================================

The SVD of a matrix $A$ is built from the symmetric matrices $A^TA$ and $AA^T$. The Spectral Theorem guarantees that these matrices have orthonormal eigenbases, which is exactly what makes the singular value decomposition possible.
